% Options for packages loaded elsewhere
\PassOptionsToPackage{unicode}{hyperref}
\PassOptionsToPackage{hyphens}{url}
\documentclass[
]{article}
\usepackage{xcolor}
\usepackage[margin=1in]{geometry}
\usepackage{amsmath,amssymb}
\setcounter{secnumdepth}{-\maxdimen} % remove section numbering
\usepackage{iftex}
\ifPDFTeX
  \usepackage[T1]{fontenc}
  \usepackage[utf8]{inputenc}
  \usepackage{textcomp} % provide euro and other symbols
\else % if luatex or xetex
  \usepackage{unicode-math} % this also loads fontspec
  \defaultfontfeatures{Scale=MatchLowercase}
  \defaultfontfeatures[\rmfamily]{Ligatures=TeX,Scale=1}
\fi
\usepackage{lmodern}
\ifPDFTeX\else
  % xetex/luatex font selection
\fi
% Use upquote if available, for straight quotes in verbatim environments
\IfFileExists{upquote.sty}{\usepackage{upquote}}{}
\IfFileExists{microtype.sty}{% use microtype if available
  \usepackage[]{microtype}
  \UseMicrotypeSet[protrusion]{basicmath} % disable protrusion for tt fonts
}{}
\makeatletter
\@ifundefined{KOMAClassName}{% if non-KOMA class
  \IfFileExists{parskip.sty}{%
    \usepackage{parskip}
  }{% else
    \setlength{\parindent}{0pt}
    \setlength{\parskip}{6pt plus 2pt minus 1pt}}
}{% if KOMA class
  \KOMAoptions{parskip=half}}
\makeatother
\usepackage{graphicx}
\makeatletter
\newsavebox\pandoc@box
\newcommand*\pandocbounded[1]{% scales image to fit in text height/width
  \sbox\pandoc@box{#1}%
  \Gscale@div\@tempa{\textheight}{\dimexpr\ht\pandoc@box+\dp\pandoc@box\relax}%
  \Gscale@div\@tempb{\linewidth}{\wd\pandoc@box}%
  \ifdim\@tempb\p@<\@tempa\p@\let\@tempa\@tempb\fi% select the smaller of both
  \ifdim\@tempa\p@<\p@\scalebox{\@tempa}{\usebox\pandoc@box}%
  \else\usebox{\pandoc@box}%
  \fi%
}
% Set default figure placement to htbp
\def\fps@figure{htbp}
\makeatother
\setlength{\emergencystretch}{3em} % prevent overfull lines
\providecommand{\tightlist}{%
  \setlength{\itemsep}{0pt}\setlength{\parskip}{0pt}}
\usepackage{bookmark}
\IfFileExists{xurl.sty}{\usepackage{xurl}}{} % add URL line breaks if available
\urlstyle{same}
\hypersetup{
  hidelinks,
  pdfcreator={LaTeX via pandoc}}

\author{}
\date{\vspace{-2.5em}}

\begin{document}

\section{Guide Animateur - Jeu des Biais
Cognitifs}\label{guide-animateur---jeu-des-biais-cognitifs}

\subsection{Formation d'une journée
complète}\label{formation-dune-journuxe9e-compluxe8te}

\begin{center}\rule{0.5\linewidth}{0.5pt}\end{center}

\subsection{📋 Vue d'ensemble}\label{vue-densemble}

\textbf{Public cible :} Débutants en biais cognitifs\\
\textbf{Durée totale :} 1 journée (7 heures)\\
\textbf{Nombre de participants :} 8 à 24 personnes\\
\textbf{Matériel nécessaire :}

\begin{itemize}
\tightlist
\item
  Jeu de 52 cartes de biais cognitifs (imprimé ou format digital)
\item
  Paperboard ou tableau
\item
  Post-its et marqueurs
\item
  Chronomètre
\item
  Badges de couleurs pour les équipes
\item
  Tableau de score
\end{itemize}

\begin{center}\rule{0.5\linewidth}{0.5pt}\end{center}

\subsection{🎯 Objectifs pédagogiques}\label{objectifs-puxe9dagogiques}

À la fin de la formation, les participants seront capables de :

\begin{enumerate}
\def\labelenumi{\arabic{enumi}.}
\tightlist
\item
  \textbf{Identifier} les 5 catégories principales de biais cognitifs
\item
  \textbf{Reconnaître} au moins 15 biais cognitifs dans des situations
  concrètes
\item
  \textbf{Analyser} l'impact des biais sur la prise de décision
  professionnelle
\item
  \textbf{Appliquer} des stratégies pour limiter l'influence des biais
  dans leur travail
\end{enumerate}

\begin{center}\rule{0.5\linewidth}{0.5pt}\end{center}

\subsection{📅 Déroulement de la
journée}\label{duxe9roulement-de-la-journuxe9e}

\subsubsection{\texorpdfstring{\textbf{Accueil et Introduction (30
min)}}{Accueil et Introduction (30 min)}}\label{accueil-et-introduction-30-min}

\paragraph{Objectifs}\label{objectifs}

\begin{itemize}
\tightlist
\item
  Créer un climat de confiance
\item
  Présenter le programme
\item
  Évaluer le niveau initial des participants
\end{itemize}

\paragraph{Déroulé détaillé}\label{duxe9rouluxe9-duxe9tailluxe9}

\textbf{1. Ice-breaker - ``La décision rapide'' (10 min)}

\begin{itemize}
\tightlist
\item
  Posez ces questions aux participants :

  \begin{itemize}
  \tightlist
  \item
    \emph{``Combien de fois levez-vous le bras par jour ?''}
  \item
    \emph{``Quel est le nombre de fenêtres dans votre ville ?''}
  \end{itemize}
\item
  Demandez-leur d'écrire leur estimation
\item
  \textbf{Révélation} : Personne ne sait ! Mais tout le monde a donné un
  chiffre → Introduction au concept d'\textbf{ancrage} et
  d'\textbf{estimation intuitive}
\end{itemize}

\textbf{2. Sondage initial (10 min)}

\begin{itemize}
\tightlist
\item
  Demandez : \emph{``Qui pense être objectif dans ses décisions ?''}
\item
  Notez le pourcentage
\item
  \textbf{Révélation} : Introduction au \textbf{biais de l'angle mort}
  (nous pensons tous être moins biaisés que les autres)
\end{itemize}

\textbf{3. Présentation du programme (10 min)}

\begin{itemize}
\tightlist
\item
  Expliquer la structure des 5 catégories de biais
\item
  Présenter les règles générales du jeu
\item
  Former 4 à 6 équipes de 3-4 personnes
\end{itemize}

\begin{quote}
{[}!IMPORTANT{]} \textbf{Point clé pour l'animateur} : Insistez dès le
départ sur le fait que TOUT LE MONDE est sujet aux biais cognitifs, y
compris vous-même. Ce n'est pas une faiblesse, mais un fonctionnement
normal du cerveau humain.
\end{quote}

\begin{center}\rule{0.5\linewidth}{0.5pt}\end{center}

\subsubsection{\texorpdfstring{\textbf{Phase 1 : Découverte des
Catégories
(1h15)}}{Phase 1 : Découverte des Catégories (1h15)}}\label{phase-1-duxe9couverte-des-catuxe9gories-1h15}

\paragraph{Objectifs}\label{objectifs-1}

\begin{itemize}
\tightlist
\item
  Découvrir les 5 catégories de biais
\item
  Comprendre la différence entre chaque catégorie
\item
  Mémoriser les symboles et couleurs
\end{itemize}

\paragraph{Distribution des cartes (15
min)}\label{distribution-des-cartes-15-min}

\textbf{Matériel} : Distribuer 10-12 cartes aléatoires par équipe

\textbf{Consigne} : \textgreater{} ``Vous avez 10 minutes pour lire vos
cartes et les classer par catégorie. Identifiez la couleur et le symbole
de chaque catégorie.''

\textbf{Les 5 catégories à identifier :}

\begin{enumerate}
\def\labelenumi{\arabic{enumi}.}
\tightlist
\item
  🧠 \textbf{Mémoire \& Souvenirs} (Violet)
\item
  🤔 \textbf{Raisonnement \& Résolution de problèmes} (Bleu)
\item
  ⚡ \textbf{Prise de décision \& Comportements} (Vert)
\item
  👥 \textbf{Travail d'équipe \& Réunions} (Orange)
\item
  🔍 \textbf{Interviews \& Tests utilisateur} (Rose)
\end{enumerate}

\paragraph{Mise en commun (20 min)}\label{mise-en-commun-20-min}

Chaque équipe présente :

\begin{itemize}
\tightlist
\item
  Les catégories identifiées
\item
  2-3 biais découverts avec leurs propres mots
\end{itemize}

\textbf{Rôle de l'animateur} :

\begin{itemize}
\tightlist
\item
  Valider ou corriger les catégories
\item
  Clarifier les définitions si nécessaire
\item
  Donner 1 exemple concret pour chaque catégorie
\end{itemize}

\paragraph{Quiz rapide - ``Quelle catégorie ?'' (15
min)}\label{quiz-rapide---quelle-catuxe9gorie-15-min}

Annoncez des situations, les équipes doivent identifier la catégorie :

\textbf{Exemples :}

\begin{enumerate}
\def\labelenumi{\arabic{enumi}.}
\tightlist
\item
  \emph{``Je me souviens mieux des images que des textes''} → 🧠 Mémoire
  (Effet de supériorité de l'image)
\item
  \emph{``On fait tous comme ça dans l'équipe''} → 👥 Travail d'équipe
  (Effet de groupe)
\item
  \emph{``Je préfère l'option par défaut''} → ⚡ Prise de décision
  (Effet par défaut)
\item
  \emph{``Si c'est nouveau, c'est forcément mieux''} → 🤔 Raisonnement
  (Biais pro-innovation)
\item
  \emph{``Je dis `oui' pour être poli''} → 🔍 Interviews (Effet de
  courtoisie)
\end{enumerate}

\textbf{Barème} : 2 points par bonne réponse

\paragraph{Synthèse théorique (25
min)}\label{synthuxe8se-thuxe9orique-25-min}

\textbf{Présentation PowerPoint ou paperboard} :

\textbf{Slide 1 : Qu'est-ce qu'un biais cognitif ?}

\begin{itemize}
\tightlist
\item
  Raccourci mental que prend notre cerveau
\item
  Économie d'énergie cognitive
\item
  Peut mener à des erreurs de jugement
\end{itemize}

\textbf{Slide 2 : Les 5 familles}

\begin{itemize}
\tightlist
\item
  Explication approfondie de chaque catégorie
\item
  Exemples dans le monde professionnel
\end{itemize}

\textbf{Slide 3 : Pourquoi c'est important ?}

\begin{itemize}
\tightlist
\item
  Impact sur la qualité des décisions
\item
  Conséquences en entreprise (recrutement, gestion de projet,
  innovation)
\end{itemize}

\begin{quote}
{[}!TIP{]} \textbf{Astuce animateur} : Préparez des exemples tirés du
secteur d'activité des participants pour maximiser l'engagement.
\end{quote}

\begin{center}\rule{0.5\linewidth}{0.5pt}\end{center}

\subsubsection{\texorpdfstring{\textbf{10h45 - 11h00 \textbar{} Pause
(15
min)}}{10h45 - 11h00 \textbar{} Pause (15 min)}}\label{h45---11h00-pause-15-min}

\begin{center}\rule{0.5\linewidth}{0.5pt}\end{center}

\subsubsection{\texorpdfstring{\textbf{11h00 - 12h30 \textbar{} Phase 2
: Jeu ``Chasseurs de biais''
(1h30)}}{11h00 - 12h30 \textbar{} Phase 2 : Jeu ``Chasseurs de biais'' (1h30)}}\label{h00---12h30-phase-2-jeu-chasseurs-de-biais-1h30}

\paragraph{Objectifs}\label{objectifs-2}

\begin{itemize}
\tightlist
\item
  Identifier des biais dans des situations réelles
\item
  Développer l'esprit critique
\item
  Apprendre par l'exemple
\end{itemize}

\paragraph{Règles du jeu (10 min)}\label{ruxe8gles-du-jeu-10-min}

\textbf{Matériel} : 15 scénarios professionnels pré-rédigés (voir Annexe
A)

\textbf{Déroulement} :

\begin{enumerate}
\def\labelenumi{\arabic{enumi}.}
\tightlist
\item
  Chaque équipe reçoit 3 scénarios différents
\item
  Pour chaque scénario, identifier \textbf{tous les biais} présents
\item
  Utiliser les cartes comme aide-mémoire
\item
  Préparer une présentation de 2 minutes
\end{enumerate}

\textbf{Barème de points} :

\begin{itemize}
\tightlist
\item
  3 points par biais correctement identifié
\item
  2 points bonus si tous les biais du scénario sont trouvés
\item
  1 point bonus pour un exemple concret similaire dans leur entreprise
\end{itemize}

\paragraph{Phase de recherche (40 min)}\label{phase-de-recherche-40-min}

\textbf{Consigne aux équipes} : \textgreater{} ``Vous avez 40 minutes
pour analyser vos 3 scénarios. Utilisez vos cartes, discutez en équipe,
et préparez votre argumentation.''

\textbf{Rôle de l'animateur pendant ce temps} :

\begin{itemize}
\tightlist
\item
  Circuler entre les équipes
\item
  Répondre aux questions
\item
  Observer les dynamiques de groupe (vous verrez les biais en action !)
\item
  Prendre des notes pour le débriefing
\end{itemize}

\begin{quote}
{[}!CAUTION{]} \textbf{Attention} : Certaines équipes peuvent tomber
dans le \textbf{biais de confirmation} (chercher seulement les biais
qu'ils connaissent déjà). Encouragez-les à explorer toutes leurs cartes.
\end{quote}

\paragraph{Présentations et débat (40
min)}\label{pruxe9sentations-et-duxe9bat-40-min}

Chaque équipe présente 1 scénario (6-7 minutes par équipe) :

\begin{itemize}
\tightlist
\item
  Lecture du scénario
\item
  Biais identifiés avec justification
\item
  Exemple personnel si disponible
\end{itemize}

\textbf{Après chaque présentation} :

\begin{itemize}
\tightlist
\item
  Les autres équipes peuvent challenger ou compléter (1 point bonus)
\item
  L'animateur valide et ajoute les biais manqués
\item
  Discussion collective : ``Comment éviter ce biais ?''
\end{itemize}

\textbf{Mise à jour du score} en direct sur le tableau

\begin{center}\rule{0.5\linewidth}{0.5pt}\end{center}

\subsubsection{\texorpdfstring{\textbf{12h30 - 13h30 \textbar{} Déjeuner
(1h)}}{12h30 - 13h30 \textbar{} Déjeuner (1h)}}\label{h30---13h30-duxe9jeuner-1h}

\begin{center}\rule{0.5\linewidth}{0.5pt}\end{center}

\subsubsection{\texorpdfstring{\textbf{13h30 - 15h00 \textbar{} Phase 3
: ``Dark UX - Le projet diabolique''
(1h30)}}{13h30 - 15h00 \textbar{} Phase 3 : ``Dark UX - Le projet diabolique'' (1h30)}}\label{h30---15h00-phase-3-dark-ux---le-projet-diabolique-1h30}

\paragraph{Objectifs}\label{objectifs-3}

\begin{itemize}
\tightlist
\item
  Comprendre la manipulation par les biais
\item
  Développer la créativité
\item
  Renforcer la mémorisation par la pratique active
\end{itemize}

\paragraph{Introduction (10 min)}\label{introduction-10-min}

\textbf{Explication du concept} : \textgreater{} ``Vous allez maintenant
utiliser vos connaissances des biais pour créer\ldots{} le produit ou
service le plus manipulateur possible ! C'est un exercice pour
comprendre les mécanismes, pas pour les appliquer !''

\textbf{Exemples de Dark UX} à montrer :

\begin{itemize}
\tightlist
\item
  Sites de réservation : ``Plus que 2 chambres !'' (rareté + urgence)
\item
  Abonnements : Option premium pré-cochée (effet par défaut)
\item
  Publicités : ``9 personnes sur 10 recommandent'' (effet de groupe +
  biais d'autorité)
\end{itemize}

\paragraph{Phase de création (50
min)}\label{phase-de-cruxe9ation-50-min}

\textbf{Consigne} : \textgreater{} ``Créez un produit, service,
publicité ou expérience en utilisant AU MINIMUM 7 biais cognitifs
différents. Vous pouvez dessiner, faire un storyboard, ou présenter un
pitch.''

\textbf{Contraintes} :

\begin{itemize}
\tightlist
\item
  Utiliser au moins 1 biais de chaque grande catégorie
\item
  Documenter précisément quel biais est utilisé et comment
\item
  Préparer une présentation de 5 minutes
\end{itemize}

\textbf{Supports fournis} :

\begin{itemize}
\tightlist
\item
  Feuilles A3, marqueurs, post-its
\item
  Accès aux 52 cartes de biais
\item
  Liste complète des biais (pour inspiration)
\end{itemize}

\textbf{Rôle de l'animateur} :

\begin{itemize}
\tightlist
\item
  Encourager la créativité
\item
  Veiller à ce que les équipes documentent bien leurs choix
\item
  Relancer celles qui bloquent avec des suggestions
\end{itemize}

\begin{quote}
{[}!WARNING{]} \textbf{Point éthique important} : Rappelez régulièrement
que c'est un exercice pédagogique. L'objectif est de comprendre pour
mieux se protéger et protéger les utilisateurs, pas de devenir
manipulateur.
\end{quote}

\paragraph{Présentations ``diaboliques'' (30
min)}\label{pruxe9sentations-diaboliques-30-min}

Chaque équipe présente son projet :

\begin{itemize}
\tightlist
\item
  Description du produit/service
\item
  Démonstration de chaque biais utilisé
\item
  Impact attendu sur l'utilisateur
\end{itemize}

\textbf{Scoring} :

\begin{itemize}
\tightlist
\item
  2 points par biais correctement appliqué
\item
  3 points bonus pour l'équipe la plus ``diabolique'' (vote des
  participants)
\item
  3 points bonus pour l'équipe la plus créative (vote de l'animateur)
\end{itemize}

\textbf{Débriefing collectif} (10 min restantes) :

\begin{itemize}
\tightlist
\item
  Comment se sentent-ils après cet exercice ?
\item
  Ont-ils reconnu des pratiques existantes ?
\item
  Comment se protéger en tant que consommateur/professionnel ?
\end{itemize}

\begin{center}\rule{0.5\linewidth}{0.5pt}\end{center}

\subsubsection{\texorpdfstring{\textbf{15h00 - 15h15 \textbar{} Pause
(15
min)}}{15h00 - 15h15 \textbar{} Pause (15 min)}}\label{h00---15h15-pause-15-min}

\begin{center}\rule{0.5\linewidth}{0.5pt}\end{center}

\subsubsection{\texorpdfstring{\textbf{15h15 - 16h30 \textbar{} Phase 4
: ``Détox cognitive'' - Stratégies anti-biais
(1h15)}}{15h15 - 16h30 \textbar{} Phase 4 : ``Détox cognitive'' - Stratégies anti-biais (1h15)}}\label{h15---16h30-phase-4-duxe9tox-cognitive---stratuxe9gies-anti-biais-1h15}

\paragraph{Objectifs}\label{objectifs-4}

\begin{itemize}
\tightlist
\item
  Apprendre à identifier ses propres biais
\item
  Découvrir des techniques de mitigation
\item
  Créer un plan d'action personnel
\end{itemize}

\paragraph{Atelier - ``Mon profil cognitif'' (25
min)}\label{atelier---mon-profil-cognitif-25-min}

\textbf{Matériel} : Questionnaire individuel (voir Annexe B)

\textbf{Consigne} : \textgreater{} ``Répondez individuellement à ces 20
situations professionnelles. Identifiez vos biais dominants.''

\textbf{Questions types} :

\begin{enumerate}
\def\labelenumi{\arabic{enumi}.}
\tightlist
\item
  \emph{Lors d'un recrutement, vous êtes influencé par :}

  \begin{itemize}
  \tightlist
  \item
    \begin{enumerate}
    \def\labelenumii{\alph{enumii})}
    \tightlist
    \item
      La première impression (ancrage)
    \end{enumerate}
  \item
    \begin{enumerate}
    \def\labelenumii{\alph{enumii})}
    \setcounter{enumii}{1}
    \tightlist
    \item
      La ressemblance avec vous (biais de similarité)
    \end{enumerate}
  \item
    \begin{enumerate}
    \def\labelenumii{\alph{enumii})}
    \setcounter{enumii}{2}
    \tightlist
    \item
      Le prestige de l'école (effet de halo)
    \end{enumerate}
  \end{itemize}
\end{enumerate}

\textbf{Résultats} : Chacun identifie ses 3-5 biais les plus fréquents

\paragraph{Présentation - ``La boîte à outils anti-biais'' (30
min)}\label{pruxe9sentation---la-bouxeete-uxe0-outils-anti-biais-30-min}

\textbf{L'animateur présente les stratégies principales} :

\textbf{1. Ralentir la décision}

\begin{itemize}
\tightlist
\item
  Prendre le temps de réfléchir
\item
  Éviter les décisions sous pression
\item
  Technique : Attendre 24h pour les décisions importantes
\end{itemize}

\textbf{2. Chercher la contradiction}

\begin{itemize}
\tightlist
\item
  Activement chercher des contre-arguments
\item
  Assigner un ``avocat du diable'' dans les réunions
\item
  Technique : ``Et si le contraire était vrai ?''
\end{itemize}

\textbf{3. Diversifier les sources}

\begin{itemize}
\tightlist
\item
  Ne pas se fier à une seule source d'information
\item
  Consulter des avis divergents
\item
  Technique : Règle des 3 sources
\end{itemize}

\textbf{4. Utiliser des checklists}

\begin{itemize}
\tightlist
\item
  Structurer la prise de décision
\item
  Évaluation systématique des critères
\item
  Technique : Grille de décision multicritères
\end{itemize}

\textbf{5. Décisions collectives structurées}

\begin{itemize}
\tightlist
\item
  Brainstorming silencieux (évite l'effet de groupe)
\item
  Vote anonyme
\item
  Technique : Méthode Delphi
\end{itemize}

\textbf{6. Attention aux contextes}

\begin{itemize}
\tightlist
\item
  Ne pas décider quand on a faim, soif, fatigue
\item
  Être conscient de son état émotionnel
\item
  Technique : Auto-questionnement avant décision
\end{itemize}

\textbf{7. Formation continue}

\begin{itemize}
\tightlist
\item
  Rester vigilant
\item
  Partager ses découvertes avec l'équipe
\item
  Technique : Rituel mensuel ``biais du mois''
\end{itemize}

\begin{quote}
{[}!TIP{]} \textbf{Pour l'animateur} : Illustrez chaque stratégie avec
un cas concret du monde professionnel des participants.
\end{quote}

\paragraph{Atelier en équipe - ``Plan d'action'' (20
min)}\label{atelier-en-uxe9quipe---plan-daction-20-min}

\textbf{Consigne} : \textgreater{} ``En équipe, choisissez 1 situation
professionnelle récurrente dans votre entreprise et proposez 3
stratégies concrètes pour limiter les biais.''

\textbf{Exemples de situations} :

\begin{itemize}
\tightlist
\item
  Réunions de lancement de projet
\item
  Processus de recrutement
\item
  Évaluation des performances
\item
  Choix de fournisseurs
\item
  Gestion de conflit
\end{itemize}

\textbf{Livrable} : Fiche action avec :

\begin{itemize}
\tightlist
\item
  Situation ciblée
\item
  Biais identifiés
\item
  3 stratégies concrètes
\item
  Indicateurs de succès
\end{itemize}

\begin{center}\rule{0.5\linewidth}{0.5pt}\end{center}

\subsubsection{\texorpdfstring{\textbf{16h30 - 17h00 \textbar{} Phase 5
: Grand Quiz Final \& Clôture (30
min)}}{16h30 - 17h00 \textbar{} Phase 5 : Grand Quiz Final \& Clôture (30 min)}}\label{h30---17h00-phase-5-grand-quiz-final-cluxf4ture-30-min}

\paragraph{Grand Quiz (20 min)}\label{grand-quiz-20-min}

\textbf{Format} : 20 questions rapides - mode ``Qui veut gagner des
millions''

\textbf{Types de questions} :

\begin{enumerate}
\def\labelenumi{\arabic{enumi}.}
\tightlist
\item
  \textbf{QCM} : ``L'effet IKEA, c'est\ldots{}''
\item
  \textbf{Vrai/Faux} : ``Les experts sont immunisés contre les biais''
\item
  \textbf{Catégorisation} : ``À quelle famille appartient le biais de
  confirmation ?''
\item
  \textbf{Situation} : ``Dans ce cas, quel biais est en jeu ?''
\item
  \textbf{Application} : ``Comment contrer le biais d'ancrage ?''
\end{enumerate}

\textbf{Barème} : 1 point par question, 5 points bonus pour l'équipe
gagnante

\textbf{Proclamation de l'équipe gagnante !} 🏆

\paragraph{Bilan et clôture (10 min)}\label{bilan-et-cluxf4ture-10-min}

\textbf{Tour de table express} :

\begin{itemize}
\tightlist
\item
  ``Quel est LE biais que vous retiendrez aujourd'hui ?''
\item
  ``Quelle action allez-vous mettre en place dès demain ?''
\end{itemize}

\textbf{Remise des supports} :

\begin{itemize}
\tightlist
\item
  Fiche récapitulative des 52 biais
\item
  Carte mentale des 5 catégories
\item
  Bibliographie et ressources complémentaires
\end{itemize}

\textbf{Évaluation de la formation} :

\begin{itemize}
\tightlist
\item
  Questionnaire de satisfaction (2 minutes)
\item
  Photo de groupe (optionnel)
\end{itemize}

\begin{center}\rule{0.5\linewidth}{0.5pt}\end{center}

\subsection{📊 Annexe A : Scénarios pour ``Chasseurs de
biais''}\label{annexe-a-scuxe9narios-pour-chasseurs-de-biais}

\subsubsection{Scénario 1 : Le recrutement
précipité}\label{scuxe9nario-1-le-recrutement-pruxe9cipituxe9}

\emph{``Sarah recrute un développeur. Le premier candidat arrive avec 15
minutes de retard et s'excuse en expliquant qu'il a eu un accident de
métro. Sarah trouve cela suspect. Le deuxième candidat, très bien
habillé, arrive à l'heure et mentionne qu'il a fait Polytechnique. Sarah
est impressionnée et décide de l'embaucher sans rencontrer d'autres
candidats, bien que sa maîtrise technique soit moyenne.''}

\textbf{Biais présents :}

\begin{itemize}
\tightlist
\item
  Biais de négativité (se souvient du retard)
\item
  Effet de halo (prestige de l'école)
\item
  Biais d'ancrage (première impression)
\item
  Biais de confirmation (cherche à confirmer sa première impression)
\item
  Effet de courtoisie (le candidat a peut-être menti)
\end{itemize}

\begin{center}\rule{0.5\linewidth}{0.5pt}\end{center}

\subsubsection{Scénario 2 : La réunion
projet}\label{scuxe9nario-2-la-ruxe9union-projet}

\emph{``En réunion de lancement, Marc, le chef de projet très
expérimenté, propose d'utiliser la même méthodologie que sur le projet
précédent qui a été un succès. Personne n'ose remettre en question cette
approche, même Julie qui pense qu'une méthode agile serait plus adaptée.
L'équipe approuve rapidement et passe à l'organisation.''}

\textbf{Biais présents :}

\begin{itemize}
\tightlist
\item
  Biais d'autorité (Marc est expérimenté)
\item
  Effet de groupe (personne n'ose contredire)
\item
  Loi de l'instrument (réutiliser ce qu'on connaît)
\item
  Biais de confirmation (Marc ne cherche pas d'alternative)
\item
  Réactance (Julie ne s'exprime pas)
\end{itemize}

\begin{center}\rule{0.5\linewidth}{0.5pt}\end{center}

\subsubsection{Scénario 3 : L'achat de
logiciel}\label{scuxe9nario-3-lachat-de-logiciel}

\emph{``L'entreprise doit choisir un nouveau logiciel de gestion. Le
commercial du Logiciel A propose une démo gratuite de 30 jours. Le
Logiciel B coûte 10 000€, le Logiciel C coûte 50 000€. Face à ces prix,
le Logiciel B semble raisonnable. Le commercial dit aussi : `Plus de 500
entreprises nous font confiance !' L'équipe choisit le Logiciel B sans
tester réellement les fonctionnalités dont elle a besoin.''}

\textbf{Biais présents :}

\begin{itemize}
\tightlist
\item
  Effet de simple exposition (démo gratuite = familiarité)
\item
  Biais d'ancrage (50 000€ rend 10 000€ attractif)
\item
  Effet de groupe (500 entreprises)
\item
  Biais d'information (trop d'infos tuent l'info)
\item
  Actualisation hyperbolique (gain immédiat : démo gratuite)
\end{itemize}

\begin{center}\rule{0.5\linewidth}{0.5pt}\end{center}

\subsubsection{Scénario 4 : Le test
utilisateur}\label{scuxe9nario-4-le-test-utilisateur}

\emph{``Thomas teste un nouveau site e-commerce. La designeuse, Sophie,
lui pose des questions : `Tu trouves que le bouton ``Acheter'' est bien
visible, non ?' Thomas répond `oui' par politesse, alors qu'il l'a
cherché 30 secondes. Sophie conclut que le design est parfait car Thomas
a validé.''}

\textbf{Biais présents :}

\begin{itemize}
\tightlist
\item
  Effet de courtoisie (Thomas est poli)
\item
  Biais de confirmation (Sophie cherche la validation)
\item
  Biais d'observation (Sophie influence l'expérience)
\item
  Illusion de validité (Sophie surinteprète les données)
\item
  Biais de l'angle mort (Sophie ne voit pas ses propres biais)
\end{itemize}

\begin{center}\rule{0.5\linewidth}{0.5pt}\end{center}

\subsubsection{Scénario 5 : L'innovation
bloquée}\label{scuxe9nario-5-linnovation-bloquuxe9e}

\emph{``L'équipe marketing propose d'utiliser TikTok pour toucher les
jeunes. Bernard, 55 ans, directeur commercial, répond : `On a toujours
fait de la publicité TV et ça marche. Les réseaux sociaux, c'est pas
pour notre cible.' Une collègue mentionne qu'un concurrent a doublé ses
ventes grâce à TikTok, mais Bernard maintient sa position : 'C'est un
coup de chance, ça ne durera pas.''}

\textbf{Biais présents :}

\begin{itemize}
\tightlist
\item
  Biais de statu quo (résistance au changement)
\item
  Loi de l'instrument (TV = outil familier)
\item
  Biais de confirmation (rejette les contre-exemples)
\item
  Dévaluation réactive (idée vient des ``jeunes'', pas de lui)
\item
  Malédiction de la connaissance (Bernard ne comprend pas la nouvelle
  génération)
\end{itemize}

\begin{center}\rule{0.5\linewidth}{0.5pt}\end{center}

\subsubsection{Scénario 6 : L'estimation
projet}\label{scuxe9nario-6-lestimation-projet}

\emph{``Julie, cheffe de projet, estime qu'un développement prendra 2
semaines. Son équipe, pourtant expérimentée, approuve. Après 2 semaines,
seulement 40\% est terminé. Julie dit : `On a eu des imprévus, personne
ne pouvait prévoir.' En réalité, TOUS leurs projets précédents ont
dépassé les délais initiaux.''}

\textbf{Biais présents :}

\begin{itemize}
\tightlist
\item
  Planification fallacieuse (sous-estimer les délais)
\item
  Biais d'autocomplaisance (blâmer les ``imprévus'')
\item
  Biais rétrospectif (auraient dû le prévoir)
\item
  Effet de groupe (l'équipe suit sans question)
\item
  Illusion de validité (confiance excessive en l'estimation)
\end{itemize}

\begin{center}\rule{0.5\linewidth}{0.5pt}\end{center}

\subsubsection{Scénario 7 : Le feedback
client}\label{scuxe9nario-7-le-feedback-client}

\emph{``Après le lancement d'un nouveau produit, l'équipe reçoit 3 avis
clients très négatifs sur Twitter. Paniqués, ils organisent une réunion
de crise et décident de tout modifier. Pendant ce temps, 500 clients
satisfaits ne laissent aucun commentaire.''}

\textbf{Biais présents :}

\begin{itemize}
\tightlist
\item
  Biais de négativité (focus sur le négatif)
\item
  Biais de disponibilité (les tweets sont immédiatement visibles)
\item
  Effet de bizarrerie (les retours négatifs marquent plus)
\item
  Actualisation hyperbolique (réaction immédiate)
\item
  Biais d'information (ignorer les 500 clients silencieux)
\end{itemize}

\begin{center}\rule{0.5\linewidth}{0.5pt}\end{center}

\subsubsection{Scénario 8 : La formation
obligatoire}\label{scuxe9nario-8-la-formation-obligatoire}

\emph{``L'entreprise impose une formation de 2 jours sur un nouveau
logiciel. Les employés rechignent : `Encore une formation inutile
imposée d'en haut !' Certains arrivent en retard, consultent leurs
téléphones. Pourtant, le contenu est pertinent et le formateur
excellent. À la fin, beaucoup admettent que c'était utile, mais
continuent à critiquer la direction pour `l'avoir imposée'.''}

\textbf{Biais présents :}

\begin{itemize}
\tightlist
\item
  Réactance (résistance car imposé)
\item
  Biais de statu quo (préfèrent leurs habitudes)
\item
  Dévaluation réactive (vient de la direction = mauvais)
\item
  Biais de confirmation (cherchent ce qui ne va pas)
\item
  Théorie de justification du système (critiquent pour maintenir leur
  vision)
\end{itemize}

\begin{center}\rule{0.5\linewidth}{0.5pt}\end{center}

\subsubsection{Scénario 9 : Le CV
prestigieux}\label{scuxe9nario-9-le-cv-prestigieux}

\emph{``Antoine, responsable RH, reçoit un CV d'un candidat ayant
travaillé chez Google, McKinsey et Harvard. Impressionné, il l'invite en
entretien sans vérifier les références. En entretien, le candidat donne
des réponses vagues, mais Antoine se dit : `Il a travaillé chez Google,
il doit être compétent.' Il l'embauche. Trois mois plus tard, les
performances sont médiocres.''}

\textbf{Biais présents :}

\begin{itemize}
\tightlist
\item
  Biais d'autorité (prestige des entreprises)
\item
  Effet de halo (Google = compétent partout)
\item
  Biais de confirmation (cherche ce qui valide sa première impression)
\item
  Illusion de validité (données prestigieuses = prédiction fiable)
\item
  Biais d'ancrage (première info : CV prestigieux)
\end{itemize}

\begin{center}\rule{0.5\linewidth}{0.5pt}\end{center}

\subsubsection{Scénario 10 : L'échec
collectif}\label{scuxe9nario-10-luxe9chec-collectif}

\emph{``Après un trimestre catastrophique, le comité de direction
analyse les résultats. Le PDG explique : `La crise économique nous a
pénalisés, nos concurrents aussi.' Le directeur marketing ajoute : `Les
clients sont devenus irrationnels.' La DRH : `Il y a eu beaucoup de
burn-out, c'est la faute au télétravail.' Personne ne mentionne la
stratégie commerciale hasardeuse décidée 6 mois plus tôt.''}

\textbf{Biais présents :}

\begin{itemize}
\tightlist
\item
  Biais d'autocomplaisance (causes externes pour les échecs)
\item
  Biais d'attribution de groupe (tout le groupe se protège)
\item
  Biais de l'angle mort (ne voient pas leur responsabilité)
\item
  Théorie de justification du système (défendent leurs décisions)
\item
  Stéréotype (« les clients sont irrationnels »)
\end{itemize}

\begin{center}\rule{0.5\linewidth}{0.5pt}\end{center}

\subsubsection{Scénario 11 : Le collègue
sympathique}\label{scuxe9nario-11-le-colluxe8gue-sympathique}

\emph{``Laura et Thomas candidatent pour le même poste de manager. Laura
est compétente mais réservée. Thomas est moins expérimenté mais très
sociable, il fait rire toute l'équipe et organise souvent des pots.
L'équipe vote pour Thomas comme futur manager. Six mois plus tard,
l'équipe souffre d'un manque d'organisation et de vision stratégique.''}

\textbf{Biais présents :}

\begin{itemize}
\tightlist
\item
  Effet de simple exposition (familiarité avec Thomas)
\item
  Effet pom-pom girl (Thomas paraît mieux en groupe)
\item
  Biais de disponibilité (souvenirs positifs de Thomas immédiatement
  accessibles)
\item
  Effet de halo (sympathie → compétence supposée)
\item
  Biais de statu quo (veulent garder l'ambiance actuelle)
\end{itemize}

\begin{center}\rule{0.5\linewidth}{0.5pt}\end{center}

\subsubsection{Scénario 12 : Le nouveau produit
innovant}\label{scuxe9nario-12-le-nouveau-produit-innovant}

\emph{``StartupTech lance SmartWidget, un objet connecté révolutionnaire
avec IA intégrée. Le marketing met en avant `Intelligence Artificielle',
`Innovation', `Futur'. Les investisseurs se bousculent. Personne ne
demande à quoi ça sert vraiment. Le produit se vend mal car les
consommateurs ne comprennent pas l'utilité. L'équipe est stupéfaite :
`Mais c'est tellement innovant !'\,''}

\textbf{Biais présents :}

\begin{itemize}
\tightlist
\item
  Biais pro-innovation (nouveau = forcément mieux)
\item
  Effet de vérité illusoire (répéter ``IA'' rend ça vrai)
\item
  Biais d'information (trop de buzzwords, peu de clarté)
\item
  Malédiction de la connaissance (comprennent le produit, oublient que
  les autres non)
\item
  Effet Dunning-Kruger potentiel (surestiment leur compréhension du
  marché)
\end{itemize}

\begin{center}\rule{0.5\linewidth}{0.5pt}\end{center}

\subsubsection{Scénario 13 : Le tableau Excel
complexe}\label{scuxe9nario-13-le-tableau-excel-complexe}

\emph{``Michel présente son analyse dans un tableau Excel de 150 lignes
avec 12 graphiques et 8 onglets. Il a passé 3 semaines dessus. La
direction, impressionnée par le travail, approuve sa recommandation sans
vraiment analyser les données. En réalité, l'analyse repose sur une
hypothèse de départ erronée, mais personne ne l'a remarqué.''}

\textbf{Biais présents :}

\begin{itemize}
\tightlist
\item
  Effet IKEA (Michel valorise son effort)
\item
  Biais d'autorité (travail complexe = compétent)
\item
  Biais d'information (trop de données tue l'analyse)
\item
  Biais de confirmation (Michel cherche ce qui valide son hypothèse)
\item
  Biais d'automatisation (confiance excessive dans Excel)
\end{itemize}

\begin{center}\rule{0.5\linewidth}{0.5pt}\end{center}

\subsubsection{Scénario 14 : L'annonce
surprise}\label{scuxe9nario-14-lannonce-surprise}

\emph{``Lors d'une réunion mensuelle habituelle, le directeur annonce
soudainement : `Nous sommes rachetés par le groupe XYZ !' Les employés
sont choqués. Trois mois plus tard, lors d'une enquête, 80\% des
employés disent : `On le savait, les signes étaient évidents.' Pourtant,
avant l'annonce, personne n'en parlait.''}

\textbf{Biais présents :}

\begin{itemize}
\tightlist
\item
  Biais rétrospectif (``je le savais'')
\item
  Illusion de validité (réinterprétation des signes passés)
\item
  Biais de confirmation rétrospectif (reconstruction de la mémoire)
\item
  Effet verbatim (se souviennent de ``l'essence'' déformée)
\end{itemize}

\begin{center}\rule{0.5\linewidth}{0.5pt}\end{center}

\subsubsection{Scénario 15 : La présentation du
lundi}\label{scuxe9nario-15-la-pruxe9sentation-du-lundi}

\emph{``Pierre présente des résultats lundi matin. Il commence par 15
minutes de contexte ennuyeux, puis présente les chiffres clés, et
termine par une conclusion spectaculaire avec une vidéo inspirante. Une
semaine plus tard, l'équipe se souvient principalement de la vidéo et de
la conclusion, mais a oublié la majorité des chiffres présentés au
milieu.''}

\textbf{Biais présents :}

\begin{itemize}
\tightlist
\item
  Effet de primauté et de récence (début et fin mieux retenus)
\item
  Effet de supériorité de l'image (vidéo \textgreater{} chiffres)
\item
  Règle du pic et de la fin (souvenir basé sur le moment fort)
\item
  Effet de bizarrerie (vidéo = inhabituel)
\item
  Effet d'humour potentiel (si vidéo drôle)
\end{itemize}

\begin{center}\rule{0.5\linewidth}{0.5pt}\end{center}

\subsection{📊 Annexe B : Questionnaire ``Mon profil
cognitif''}\label{annexe-b-questionnaire-mon-profil-cognitif}

\textbf{Instructions} : Pour chaque situation, cochez la réponse qui
vous correspond le PLUS souvent. Soyez honnête !

\subsubsection{Section 1 : Prise de
décision}\label{section-1-prise-de-duxe9cision}

\textbf{1. Lors d'un achat important, vous :}

\begin{itemize}
\tightlist
\item
  \begin{enumerate}
  \def\labelenumi{\alph{enumi})}
  \tightlist
  \item
    Choisissez rapidement ce qui vous plaît au premier coup d'œil
    \emph{(Ancrage)}
  \end{enumerate}
\item
  \begin{enumerate}
  \def\labelenumi{\alph{enumi})}
  \setcounter{enumi}{1}
  \tightlist
  \item
    Prenez l'option recommandée par le vendeur \emph{(Autorité)}
  \end{enumerate}
\item
  \begin{enumerate}
  \def\labelenumi{\alph{enumi})}
  \setcounter{enumi}{2}
  \tightlist
  \item
    Choisissez l'option par défaut ou moyenne \emph{(Défaut / Ancrage)}
  \end{enumerate}
\item
  \begin{enumerate}
  \def\labelenumi{\alph{enumi})}
  \setcounter{enumi}{3}
  \tightlist
  \item
    Comparez systématiquement toutes les options \emph{(Stratégie
    rationnelle)}
  \end{enumerate}
\end{itemize}

\textbf{2. Face à un nouveau projet, vous :}

\begin{itemize}
\tightlist
\item
  \begin{enumerate}
  \def\labelenumi{\alph{enumi})}
  \tightlist
  \item
    Réutilisez les méthodes qui ont déjà fonctionné \emph{(Loi de
    l'instrument)}
  \end{enumerate}
\item
  \begin{enumerate}
  \def\labelenumi{\alph{enumi})}
  \setcounter{enumi}{1}
  \tightlist
  \item
    Cherchez ce que font les autres entreprises \emph{(Effet de groupe)}
  \end{enumerate}
\item
  \begin{enumerate}
  \def\labelenumi{\alph{enumi})}
  \setcounter{enumi}{2}
  \tightlist
  \item
    Privilégiez les solutions innovantes même risquées \emph{(Biais
    pro-innovation)}
  \end{enumerate}
\item
  \begin{enumerate}
  \def\labelenumi{\alph{enumi})}
  \setcounter{enumi}{3}
  \tightlist
  \item
    Analysez les besoins avant de choisir la méthode \emph{(Stratégie
    rationnelle)}
  \end{enumerate}
\end{itemize}

\textbf{3. Quand quelqu'un vous contredit, vous :}

\begin{itemize}
\tightlist
\item
  \begin{enumerate}
  \def\labelenumi{\alph{enumi})}
  \tightlist
  \item
    Cherchez des arguments pour confirmer votre position
    \emph{(Confirmation)}
  \end{enumerate}
\item
  \begin{enumerate}
  \def\labelenumi{\alph{enumi})}
  \setcounter{enumi}{1}
  \tightlist
  \item
    Défendez votre point de vue avec énergie \emph{(Réactance)}
  \end{enumerate}
\item
  \begin{enumerate}
  \def\labelenumi{\alph{enumi})}
  \setcounter{enumi}{2}
  \tightlist
  \item
    Accordez plus de poids à vos arguments qu'aux leurs
    \emph{(Autocomplaisance)}
  \end{enumerate}
\item
  \begin{enumerate}
  \def\labelenumi{\alph{enumi})}
  \setcounter{enumi}{3}
  \tightlist
  \item
    Écoutez et évaluez objectivement \emph{(Stratégie rationnelle)}
  \end{enumerate}
\end{itemize}

\subsubsection{Section 2 : Travail
d'équipe}\label{section-2-travail-duxe9quipe}

\textbf{4. En réunion, vous :}

\begin{itemize}
\tightlist
\item
  \begin{enumerate}
  \def\labelenumi{\alph{enumi})}
  \tightlist
  \item
    Suivez l'avis majoritaire pour éviter les conflits \emph{(Effet de
    groupe)}
  \end{enumerate}
\item
  \begin{enumerate}
  \def\labelenumi{\alph{enumi})}
  \setcounter{enumi}{1}
  \tightlist
  \item
    Ne partagez pas votre désaccord pour rester poli \emph{(Courtoisie)}
  \end{enumerate}
\item
  \begin{enumerate}
  \def\labelenumi{\alph{enumi})}
  \setcounter{enumi}{2}
  \tightlist
  \item
    Pensez que votre expertise vous rend plus objectif que les autres
    \emph{(Angle mort)}
  \end{enumerate}
\item
  \begin{enumerate}
  \def\labelenumi{\alph{enumi})}
  \setcounter{enumi}{3}
  \tightlist
  \item
    Exprimez votre position même si elle diverge \emph{(Stratégie
    rationnelle)}
  \end{enumerate}
\end{itemize}

\textbf{5. Face à une idée d'un concurrent, vous :}

\begin{itemize}
\tightlist
\item
  \begin{enumerate}
  \def\labelenumi{\alph{enumi})}
  \tightlist
  \item
    La critiquez automatiquement \emph{(Dévaluation réactive)}
  \end{enumerate}
\item
  \begin{enumerate}
  \def\labelenumi{\alph{enumi})}
  \setcounter{enumi}{1}
  \tightlist
  \item
    Préférez les idées développées en interne \emph{(NIH)}
  \end{enumerate}
\item
  \begin{enumerate}
  \def\labelenumi{\alph{enumi})}
  \setcounter{enumi}{2}
  \tightlist
  \item
    L'évaluez selon qui la propose \emph{(Autorité)}
  \end{enumerate}
\item
  \begin{enumerate}
  \def\labelenumi{\alph{enumi})}
  \setcounter{enumi}{3}
  \tightlist
  \item
    L'analysez sur ses mérites propres \emph{(Stratégie rationnelle)}
  \end{enumerate}
\end{itemize}

\textbf{6. Lorsque votre équipe réussit, c'est grâce à :}

\begin{itemize}
\tightlist
\item
  \begin{enumerate}
  \def\labelenumi{\alph{enumi})}
  \tightlist
  \item
    Vos efforts et compétences \emph{(Autocomplaisance)}
  \end{enumerate}
\item
  \begin{enumerate}
  \def\labelenumi{\alph{enumi})}
  \setcounter{enumi}{1}
  \tightlist
  \item
    La synergie d'équipe \emph{(Attribution groupe)}
  \end{enumerate}
\item
  \begin{enumerate}
  \def\labelenumi{\alph{enumi})}
  \setcounter{enumi}{2}
  \tightlist
  \item
    Les bonnes méthodes que vous utilisez \emph{(Justification système)}
  \end{enumerate}
\item
  \begin{enumerate}
  \def\labelenumi{\alph{enumi})}
  \setcounter{enumi}{3}
  \tightlist
  \item
    Un mix de compétences, efforts et circonstances \emph{(Vision
    équilibrée)}
  \end{enumerate}
\end{itemize}

\subsubsection{Section 3 : Mémoire et
information}\label{section-3-muxe9moire-et-information}

\textbf{7. Vous retenez surtout :}

\begin{itemize}
\tightlist
\item
  \begin{enumerate}
  \def\labelenumi{\alph{enumi})}
  \tightlist
  \item
    Les informations récentes et disponibles \emph{(Disponibilité)}
  \end{enumerate}
\item
  \begin{enumerate}
  \def\labelenumi{\alph{enumi})}
  \setcounter{enumi}{1}
  \tightlist
  \item
    Ce qui confirme vos croyances \emph{(Confirmation)}
  \end{enumerate}
\item
  \begin{enumerate}
  \def\labelenumi{\alph{enumi})}
  \setcounter{enumi}{2}
  \tightlist
  \item
    Les événements marquants et inhabituels \emph{(Bizarrerie)}
  \end{enumerate}
\item
  \begin{enumerate}
  \def\labelenumi{\alph{enumi})}
  \setcounter{enumi}{3}
  \tightlist
  \item
    Vous prenez des notes pour tout \emph{(Stratégie rationnelle)}
  \end{enumerate}
\end{itemize}

\textbf{8. Face à un problème, vous cherchez des informations :}

\begin{itemize}
\tightlist
\item
  \begin{enumerate}
  \def\labelenumi{\alph{enumi})}
  \tightlist
  \item
    Jusqu'à être certain d'avoir tout \emph{(Biais d'information)}
  \end{enumerate}
\item
  \begin{enumerate}
  \def\labelenumi{\alph{enumi})}
  \setcounter{enumi}{1}
  \tightlist
  \item
    Seulement ce qui confirme votre intuition \emph{(Confirmation)}
  \end{enumerate}
\item
  \begin{enumerate}
  \def\labelenumi{\alph{enumi})}
  \setcounter{enumi}{2}
  \tightlist
  \item
    Sur Google, c'est plus rapide \emph{(Effet Google)}
  \end{enumerate}
\item
  \begin{enumerate}
  \def\labelenumi{\alph{enumi})}
  \setcounter{enumi}{3}
  \tightlist
  \item
    Suffisamment pour décider, pas plus \emph{(Stratégie rationnelle)}
  \end{enumerate}
\end{itemize}

\subsubsection{Section 4 : Évaluation des
autres}\label{section-4-uxe9valuation-des-autres}

\textbf{9. Lors d'un recrutement, vous êtes influencé par :}

\begin{itemize}
\tightlist
\item
  \begin{enumerate}
  \def\labelenumi{\alph{enumi})}
  \tightlist
  \item
    La première impression \emph{(Ancrage / Primauté)}
  \end{enumerate}
\item
  \begin{enumerate}
  \def\labelenumi{\alph{enumi})}
  \setcounter{enumi}{1}
  \tightlist
  \item
    Le prestige de la formation \emph{(Autorité / Halo)}
  \end{enumerate}
\item
  \begin{enumerate}
  \def\labelenumi{\alph{enumi})}
  \setcounter{enumi}{2}
  \tightlist
  \item
    La ressemblance avec vous \emph{(Similarité)}
  \end{enumerate}
\item
  \begin{enumerate}
  \def\labelenumi{\alph{enumi})}
  \setcounter{enumi}{3}
  \tightlist
  \item
    Uniquement les compétences vérifiées \emph{(Stratégie rationnelle)}
  \end{enumerate}
\end{itemize}

\textbf{10. Vous évaluez la performance d'un collègue sur :}

\begin{itemize}
\tightlist
\item
  \begin{enumerate}
  \def\labelenumi{\alph{enumi})}
  \tightlist
  \item
    Ses dernières actions \emph{(Disponibilité / Récence)}
  \end{enumerate}
\item
  \begin{enumerate}
  \def\labelenumi{\alph{enumi})}
  \setcounter{enumi}{1}
  \tightlist
  \item
    Votre relation personnelle avec lui \emph{(Halo)}
  \end{enumerate}
\item
  \begin{enumerate}
  \def\labelenumi{\alph{enumi})}
  \setcounter{enumi}{2}
  \tightlist
  \item
    Les moments forts, bons ou mauvais \emph{(Pic et fin)}
  \end{enumerate}
\item
  \begin{enumerate}
  \def\labelenumi{\alph{enumi})}
  \setcounter{enumi}{3}
  \tightlist
  \item
    Une évaluation continue sur toute l'année \emph{(Stratégie
    rationnelle)}
  \end{enumerate}
\end{itemize}

\subsubsection{Interprétation}\label{interpruxe9tation}

\textbf{Comptez vos réponses a, b, c, d} :

\begin{itemize}
\item
  \textbf{Majorité de a, b ou c} : Vous êtes comme tout le monde ! Les
  biais sont actifs. Identifiez vos patterns (a = ancrage, b =
  autorité/groupe, c = confirmation/innovation).
\item
  \textbf{Majorité de d} : Soit vous êtes très conscient des biais et
  appliquez déjà des stratégies rationnelles, soit vous êtes victime du
  \textbf{biais d'angle mort} (vous pensez être plus rationnel que vous
  ne l'êtes réellement) ! 😊
\end{itemize}

\textbf{Profil dominant} :

\begin{itemize}
\tightlist
\item
  \textbf{Si beaucoup de ``a''} → Tendance \textbf{Ancrage /
  Disponibilité} : Vous privilégiez les premières informations
\item
  \textbf{Si beaucoup de ``b''} → Tendance \textbf{Autorité / Groupe} :
  Vous êtes influencé par les autres
\item
  \textbf{Si beaucoup de ``c''} → Tendance \textbf{Confirmation /
  Innovation} : Vous cherchez ce qui valide vos croyances
\end{itemize}

\begin{center}\rule{0.5\linewidth}{0.5pt}\end{center}

\subsection{📚 Annexe C : Ressources
complémentaires}\label{annexe-c-ressources-compluxe9mentaires}

\subsubsection{Livres recommandés}\label{livres-recommanduxe9s}

\begin{itemize}
\tightlist
\item
  \textbf{``Thinking, Fast and Slow''} - Daniel Kahneman
\item
  \textbf{``Nudge''} - Richard Thaler \& Cass Sunstein
\item
  \textbf{``Predictably Irrational''} - Dan Ariely
\item
  \textbf{``The Art of Thinking Clearly''} - Rolf Dobelli
\end{itemize}

\subsubsection{Sites web}\label{sites-web}

\begin{itemize}
\tightlist
\item
  \url{https://www.cognitive-biases.org} - Carte interactive
\item
  \url{https://thedecisionlab.com} - Bibliothèque de biais
\item
  Codex des biais cognitifs
\end{itemize}

\subsubsection{Vidéos pédagogiques}\label{viduxe9os-puxe9dagogiques}

\begin{itemize}
\tightlist
\item
  TED Talks : chercher ``cognitive bias'', ``Daniel Kahneman'',
  ``behavioral economics''
\item
  Chaîne YouTube : ``Veritasium'' (biais scientifiques)
\end{itemize}

\begin{center}\rule{0.5\linewidth}{0.5pt}\end{center}

\subsection{✅ Checklist de préparation pour
l'animateur}\label{checklist-de-pruxe9paration-pour-lanimateur}

\subsubsection{\texorpdfstring{\textbf{1 semaine avant
:}}{1 semaine avant :}}\label{semaine-avant}

\begin{itemize}
\tightlist
\item[$\square$]
  Imprimer ou préparer les 52 cartes de biais (1 jeu par équipe + 1 pour
  l'animateur)
\item[$\square$]
  Préparer les 15 scénarios ``Chasseurs de biais'' (3 par équipe)
\item[$\square$]
  Créer le support de présentation (slides)
\item[$\square$]
  Imprimer les questionnaires ``Mon profil cognitif''
\item[$\square$]
  Préparer le tableau de score
\item[$\square$]
  Réserver la salle et vérifier l'équipement (vidéoprojecteur,
  paperboard)
\end{itemize}

\subsubsection{\texorpdfstring{\textbf{La veille
:}}{La veille :}}\label{la-veille}

\begin{itemize}
\tightlist
\item[$\square$]
  Photocopier les documents participants (règles, fiches récap)
\item[$\square$]
  Préparer le matériel créatif (marqueurs, post-its, feuilles A3)
\item[$\square$]
  Tester le vidéoprojecteur et le chronomètre
\item[$\square$]
  Préparer les badges d'équipe (optionnel)
\item[$\square$]
  Prévoir des récompenses symboliques pour l'équipe gagnante
\end{itemize}

\subsubsection{\texorpdfstring{\textbf{Le jour J
:}}{Le jour J :}}\label{le-jour-j}

\begin{itemize}
\tightlist
\item[$\square$]
  Arriver 30 min avant pour installer
\item[$\square$]
  Mettre en place les tables en îlots (1 par équipe)
\item[$\square$]
  Disposer le matériel sur chaque table
\item[$\square$]
  Afficher le programme de la journée
\item[$\square$]
  Prévoir eau, café, collations si possible
\item[$\square$]
  Avoir votre propre jeu de cartes pour démonstration
\end{itemize}

\begin{center}\rule{0.5\linewidth}{0.5pt}\end{center}

\subsection{💡 Conseils d'animation}\label{conseils-danimation}

\subsubsection{\texorpdfstring{\textbf{Gérer les participants
sceptiques}}{Gérer les participants sceptiques}}\label{guxe9rer-les-participants-sceptiques}

Certains peuvent penser être ``immunisés'' contre les biais
(\textbf{biais de l'angle mort} !).

\begin{itemize}
\tightlist
\item
  Désamorcez avec humour : ``Si vous pensez ne pas avoir de biais, vous
  avez le biais de l'angle mort !''
\item
  Utilisez des démonstrations visuelles (illusions d'optique = biais
  visuel)
\item
  Partagez vos propres erreurs liées aux biais
\end{itemize}

\subsubsection{\texorpdfstring{\textbf{Maintenir
l'énergie}}{Maintenir l'énergie}}\label{maintenir-luxe9nergie}

\begin{itemize}
\tightlist
\item
  Variez les formats : individuel, binôme, équipe, collectif
\item
  Utilisez le mouvement : faites lever les participants pour voter
\item
  Célébrez les bonnes réponses avec enthousiasme
\item
  Utilisez des exemples drôles ou surprenants
\end{itemize}

\subsubsection{\texorpdfstring{\textbf{Gérer le
timing}}{Gérer le timing}}\label{guxe9rer-le-timing}

\begin{itemize}
\tightlist
\item
  Utilisez un chronomètre visible pour les exercices
\item
  Prévoyez 10\% de marge sur chaque activité
\item
  Si vous êtes en retard : supprimez plutôt un exemple qu'une activité
  entière
\item
  Les phases de présentation peuvent déborder : soyez flexible
\end{itemize}

\subsubsection{\texorpdfstring{\textbf{Adapter au
niveau}}{Adapter au niveau}}\label{adapter-au-niveau}

\begin{itemize}
\tightlist
\item
  \textbf{Si le groupe est très avancé} : Ajoutez des biais
  supplémentaires, proposez des cas plus complexes
\item
  \textbf{Si le groupe est en difficulté} : Simplifiez les scénarios,
  donnez plus d'indices, réduisez le nombre de biais par exercice
\end{itemize}

\subsubsection{\texorpdfstring{\textbf{Créer une dynamique
positive}}{Créer une dynamique positive}}\label{cruxe9er-une-dynamique-positive}

\begin{itemize}
\tightlist
\item
  Évitez le jugement : ``C'est normal d'avoir des biais''
\item
  Valorisez la prise de conscience : ``Bravo d'avoir identifié ce biais
  chez vous''
\item
  Encouragez le partage d'expériences personnelles
\item
  Terminez sur une note positive et actionnable
\end{itemize}

\begin{center}\rule{0.5\linewidth}{0.5pt}\end{center}

\subsection{🎯 Indicateurs de succès de la
formation}\label{indicateurs-de-succuxe8s-de-la-formation}

\subsubsection{Pendant la formation}\label{pendant-la-formation}

\begin{itemize}
\tightlist
\item
  ✅ Participation active de tous les membres
\item
  ✅ Questions pertinentes posées
\item
  ✅ Exemples personnels partagés spontanément
\item
  ✅ Débats constructifs entre participants
\item
  ✅ Rires et bonne ambiance
\end{itemize}

\subsubsection{Après la formation}\label{apruxe8s-la-formation}

\begin{itemize}
\tightlist
\item
  ✅ Score moyen au quiz final \textgreater{} 70\%
\item
  ✅ Satisfaction générale \textgreater{} 4/5
\item
  ✅ Chaque participant identifie au moins 3 actions concrètes
\item
  ✅ Demandes de ressources complémentaires
\item
  ✅ Application des concepts dans les semaines suivantes (à suivre en
  questionnaire J+30)
\end{itemize}

\begin{center}\rule{0.5\linewidth}{0.5pt}\end{center}

\textbf{Bonne animation ! 🎉}

N'oubliez pas : En tant qu'animateur, vous êtes aussi sujet aux biais
cognitifs. Restez humble, curieux, et amusez-vous !

\end{document}
